\chapter{Fundamentação teórica}
\label{cap:fundamentacao-teorica}

% Conteúdo realocado da introdução. O capítulo será ampliado na etapa de fundamentação teórica.
\section{Pré-dimensionamento de barras de palitos}
\label{sec:fund-pre-dimensionamento}

O pré-dimensionamento de pontes de palitos relaciona os esforços obtidos na análise estrutural às propriedades do material e à geometria dos elementos. Essa relação integra o processo de definição da composição das barras e da quantidade de material necessária à construção \cite{silva2024metodologia}.

Nas barras submetidas à tração, a avaliação considera a relação entre força normal, área resistente e resistência adotada para o material. Nas barras comprimidas, deve-se considerar também a possibilidade de instabilidade por flambagem, cujo estudo integra a análise de elementos de treliças \cite{souza2015flambagem}. Ao aplicar esses conceitos aos palitos, é necessário distinguir o número de camadas que compõem a seção transversal da quantidade física de palitos distribuídos ao longo da barra.

Essas estimativas dependem das propriedades da madeira, das hipóteses de cálculo e das condições de construção. Por isso, a verificação dos resultados de um programa não substitui a avaliação experimental do comportamento da ponte. A relação entre construção, propriedades do material e desempenho estrutural integra a metodologia apresentada por \citeonline{silva2024metodologia}.
